\chapter{結言}
	\section{まとめ}
% 	本論文では，ロボットにより容器に液体を定量的に注ぐための手法として，RGB-Dカメラを使用して目標容器の形状を取得し，それを用いて注がれた液体の体積を推定することによって指定した量の液体を注ぐ手法を提案した．
% 	そして，その手法を双腕ロボットに実装し，体積測定実験と注ぐ実験を行い，計3種類の目標容器に定量の液体を注いだ．
% 	実験から得られたデータで，平均$\pm4\ \mathrm{ml}$の誤差があるが，日常生活において十分な精度であると考えられる．
% 	したがって，本稿で提案した手法は容器に液体を定量的に注ぐことができる．
	
	本論文では，双腕ロボットと視覚情報を用いて不透明な液体を目標容器に定量的に注ぐための2つの手法を提案し，実装した．
	\begin{enumerate}
	\item 手法1は注ぎ先の容器に注いだ液体の液面高さに注目し，注ぎ先の容器のモデルを使用し，注ぐ時に液面高さを注いだ体積にリアルタイムに変換し，注ぐ容器の傾き角速度を操作量，注いだ体積を制御量とするPD制御で注ぐ作業を実現する手法である．
	\item 手法2は注ぐ容器モデルと注ぐ容器の初期液体体積を計測し，注ぐ容器傾き角度と注いだ体積の関係により，注ぐ容器を目標体積に対応する傾き角度に傾けることで注ぐ作業を実現する手法である．
	\end{enumerate}
	
	そして，それらの手法を双腕ロボットシステムに実装した．
	更に，実験のために，多視点点群の取得と合成や液面検出や変換テーブルの計算と連続化などのプログラムも設計し，実装した．
	
	また，共通実験である体積測定の実験を行い，手法1は計3種類の目標容器に定量の液体を注ぐ実験をそれぞれ行い，手法2は計2種類の注ぐ容器を使用し，注ぐ実験をそれぞれ行った．
	
	実験結果により，手法1（平均$\pm4\ \mathrm{ml}$で，先行研究\cite{Related1}より5倍以上の精度）は定量的に注げることを確認した．
	手法2は今回の二種類の容器において20.79[ml]の精度で定量的に注げることを確認した．
	
	したがって，本研究は注ぐ容器と注ぎ先の容器のモデルどちらかが計測できる条件で液体を目標容器に定量的に注ぐことを実現した．
	
	\section{今後の展望}
	今回の体積の測定と注ぐ縁の検出の実装において，それぞれで断面の処理が必要である．
	その処理について，今回はConvexHullを使用し，断面積の計算と断面の不連続な箇所の検出をした．
	しかし，ConvexHullを使用したため，もし断面の形状が凹んでいるなら，断面積の計算に間違いが生じ，凹みのある場所が注ぎ口と誤認されてしまった．
	そこで，断面の処理はConvexHullではなく，ConcaveHullを使用し，より複雑な形状の容器にも対応できるようにする予定である．
	
	次は，表面張力について，手法2は注いでいる時に注ぐ容器の液面の高さが注ぐポイントの高さと同じであることを前提とした．しかし，現実は容器の材質と液体の種類による表面張力の影響があるので，液面の高さと注ぐポイントの高さに差異が生じてしまい，その分の体積の差異は一部の誤差を発生させた．
	したがって，今後は表面張力を考慮する傾き角度と注いだ体積の関係を計算する方法を見つける予定である．